\documentclass[conference]{IEEEtran}
\IEEEoverridecommandlockouts

% ── Packages ──────────────────────────────────────────────────────────────────
\usepackage{cite}
\usepackage{amsmath,amssymb}
\usepackage{graphicx}
\usepackage{xcolor}
\usepackage{hyperref}
\usepackage{microtype}
\usepackage{algorithm}
\usepackage{algpseudocode}
\usepackage{booktabs}
\usepackage{enumitem}

% ── Title block ───────────────────────────────────────────────────────────────
\title{AdaptiveServe: Runtime-Adaptive LLM Inference via KV Cache Strategy
       Selection, Activation Sparsity, and Model Routing}

\author{
  \IEEEauthorblockN{Youssef Araby}
  \IEEEauthorblockA{\textit{[Department / Lab]}\\
    [Institution]\\
    [City, Country]\\
    youssef@[institution].edu}
  \and
  \IEEEauthorblockN{Rodaina [LastName]}
  \IEEEauthorblockA{\textit{[Department / Lab]}\\
    [Institution]\\
    [City, Country]\\
    rodaina@[institution].edu}
}

% ══════════════════════════════════════════════════════════════════════════════
\begin{document}

\maketitle

% ── Abstract ──────────────────────────────────────────────────────────────────
\begin{abstract}
Efficient LLM inference is bottlenecked by KV cache memory bandwidth and dense
compute, yet existing compression methods apply a single fixed policy to every
prompt.
We present \textsc{AdaptiveServe}, a five-layer adaptive inference framework;
this paper covers layers (i)--(iii).
Our KV cache selector uses three cheap prefill signals---attention entropy,
sequence length, and head-level variance---in a two-phase design: a
rule-based Phase~A heuristic and a learned Phase~B MLP.
An initial imbalanced dataset caused Ada-KV to be never selected; expanding
the calibration corpus resolved this, yielding 61\% classifier accuracy vs.\
a 49\% majority baseline under 5-fold cross-validation and a Pareto frontier
that strictly dominates every individual technique.
Integrating \textsc{TailorKV}---an offline method that degrades on some tasks
alone---into the classifier further improves the Pareto frontier by routing
away from it on quality-sensitive inputs.
For activation sparsity, Wanda 2:4 on Llama-3.1-8B is blocked by two issues:
vLLM~0.20 removed \texttt{Sparse24} kernel support (eliminating any
throughput gain) and insufficient calibration (512 samples,
\textsc{OpenPlatypus}) caused an 18~pp PIQA accuracy drop.
For model routing, Phi-3.5-Mini (3.8B) and Llama-3.1-8B-Instruct perform
comparably on LongBench and WikiText, making reliable routing signal
unlearnable; effective routing requires a larger capability gap between models.
\end{abstract}

\begin{IEEEkeywords}
LLM inference, KV cache compression, adaptive selection, weight sparsity,
model routing
\end{IEEEkeywords}

% ══════════════════════════════════════════════════════════════════════════════
%  I. INTRODUCTION
% ══════════════════════════════════════════════════════════════════════════════
\section{Introduction}
\label{sec:introduction}

The cost of LLM inference has become a first-class engineering concern.
Serving a single 7--8B parameter model at even moderate traffic imposes three
interacting constraints: memory bandwidth to stream the KV cache from DRAM to
GPU on every decode step, compute throughput to execute dense GEMMs in
attention and feed-forward layers, and a latency budget that varies by
application---a coding assistant can tolerate higher time-to-first-token than a
real-time speech interface.
These constraints do not decompose cleanly: reducing KV cache size via
quantization saves memory bandwidth but can degrade generation quality;
applying weight sparsity saves FLOPs but only if the inference engine exposes
the corresponding sparse kernels; routing short prompts to a smaller model
saves compute but only if the smaller model is genuinely weaker on those
inputs.

The literature addresses each axis in isolation.
KV cache compression methods---KVQuant~\cite{hooper2024kvquant},
DynamicKV~\cite{dynamickv}, Ada-KV~\cite{adakv}, and
TailorKV~\cite{tailorkv}---define a fixed per-layer compression policy applied
uniformly to all inputs, without considering how task type, sequence length, or
attention statistics should modulate the choice.
Weight sparsity methods---Wanda~\cite{sun2024wanda} and
SparseGPT~\cite{frantar2023sparsegpt}---reduce the density of model weights
but require hardware-friendly sparsity patterns and a serving stack that
exposes the corresponding sparse dispatch kernels.
Speculative decoding and model routing methods~\cite{leviathan2023speculative,%
chen2023speculative} exploit heterogeneity across model sizes but their gains
diminish when the draft and target models are close in capability.

\textsc{AdaptiveServe} brings these axes together under a unified five-layer
framework: (i)~adaptive KV cache strategy selection, (ii)~activation sparsity,
(iii)~query-level model routing, (iv)~cloud--edge splitting, and (v)~serving-
level batching.
This paper reports findings on layers (i)--(iii), for which we have completed
experiments using LLaMA-3.1-8B-Instruct~\cite{llama3} and
Phi-3.5-Mini-3.8B~\cite{phi3} served via vLLM~\cite{vllm}, evaluated on
LongBench~\cite{longbench}, WikiText-2, and PIQA~\cite{piqa} using the
EleutherAI evaluation harness~\cite{lmeval}.

\subsection{Motivation: Adaptive KV Cache Selection}
\label{sec:intro:kvcache}

A central observation is that no single KV cache compression policy is
Pareto-optimal across the full distribution of LLM prompts.
Short prompts with low attention entropy lose little from aggressive KV
quantization; long prompts with high head-level variance benefit from
head-adaptive budgets (\textsc{Ada-KV}~\cite{adakv}); certain task types are
sensitive to 3-bit quantization artifacts that 4-bit quantization avoids
(\textsc{KVQuant}~\cite{hooper2024kvquant}); and some prompts tolerate
extreme compression ratios that \textsc{TailorKV}~\cite{tailorkv} provides
while others do not.
A selector that dispatches to the right policy using cheap prefill-time signals
can, in principle, achieve a strictly better compression--quality Pareto
frontier than any fixed policy.

This motivates the key design question: \emph{what signals are cheap enough to
compute per layer at prefill time and informative enough to distinguish among
the four target strategies?}
We identify three such signals---attention entropy, sequence length, and
head-level variance---all of which are available during the prefill pass with
negligible additional overhead.
Even a small MLP trained on approximately 2\,800 labeled examples achieves
meaningful routing accuracy and, critically, produces a Pareto frontier that
dominates every individual compression technique.

\subsection{Motivation: Limits of Weight Sparsity}
\label{sec:intro:sparsity}

An alternative to KV cache compression is to reduce the computational cost
of the attention and feed-forward layers via weight sparsity.
R-Sparse~\cite{rsparse} claims dynamic activation sparsity with minimal
accuracy loss, but the custom CUDA kernels required to realize its reported
latency improvements are not publicly available; we were unable to reproduce
the speedup numbers.
Wanda 2:4~\cite{sun2024wanda} offers a tractable alternative with an open
implementation via \texttt{llm-compressor}, and we successfully applied it to
Llama-3.1-8B.
However, our evaluation revealed two blocking issues on the current vLLM~0.20
stack.
First, vLLM~0.20 removed the \texttt{Sparse24} kernel dispatch path
(\texttt{NotImplementedError: Sparse24 models are no longer supported by
vLLM}), meaning any Wanda~2:4 checkpoint---including those declared via a
\texttt{sparsity\_config} block---fails to load or falls back to dense
kernels with no throughput benefit.
Second, calibrating on only 512 samples from \textsc{OpenPlatypus} yielded an
18~pp PIQA accuracy drop (0.6300 vs.\ the dense baseline of 0.8118), far
outside the 1--3~pp range reported by Sun~et~al.~\cite{sun2024wanda} under
proper calibration.
We document both issues in detail in Section~\ref{sec:sparsity} as a
cautionary note for practitioners relying on planning documents that predate
vLLM~0.20.

\subsection{Motivation: Limits of Model Routing}
\label{sec:intro:routing}

Routing easy prompts to a small model and hard prompts to a large model is
appealing in theory.
In our evaluation, however, Phi-3.5-Mini-3.8B and Llama-3.1-8B-Instruct
perform comparably on LongBench and WikiText-2 perplexity.
When the capability gap between models is small, the routing signal that
consistently favors one model over the other is difficult to learn: the router
degenerates either to always using the large model (safe but uneconomical) or
to routing too aggressively to the small model (cheap but lower quality).
We conclude that effective model routing requires a pair of models with
substantially different capability profiles, and we outline the conditions
under which routing is expected to yield meaningful gains in
Section~\ref{sec:routing}.

\subsection{Paper Organization}
\label{sec:intro:organization}

Section~\ref{sec:related} reviews related work on KV cache compression, weight
sparsity, and model routing.
Section~\ref{sec:kvcache} describes the KV cache selector in full: the
Phase~A rule-based heuristic, the calibration pipeline, the Phase~B MLP
classifier, and TailorKV integration.
Section~\ref{sec:sparsity} documents our activation sparsity experiments,
including the Wanda~2:4 pipeline, the vLLM \texttt{Sparse24} removal, and the
accuracy analysis.
Section~\ref{sec:routing} presents the model routing study comparing
Phi-3.5-Mini and Llama-3.1-8B-Instruct.
Section~\ref{sec:discussion} synthesizes findings across all three axes and
identifies the conditions under which each technique yields practical benefit.
Section~\ref{sec:conclusion} concludes and outlines future work.

% ── Forward-reference stubs for sections to be added in future files ────────
% When Sections II–VII are written, move each \label{} to the corresponding
% \section{} command and remove the stub declarations below.
\label{sec:related}   % TODO: move to \section{Related Work}
\label{sec:kvcache}   % TODO: move to \section{KV Cache Selector}
\label{sec:sparsity}  % TODO: move to \section{Activation Sparsity}
\label{sec:routing}   % TODO: move to \section{Model Routing}
\label{sec:discussion}% TODO: move to \section{Discussion}
\label{sec:conclusion}% TODO: move to \section{Conclusion}

% ── Bibliography ──────────────────────────────────────────────────────────────
\bibliographystyle{IEEEtran}
\bibliography{adaptiveserve}

\end{document}
